\section{Exact current sufficiency for the final hidden layer}
The final hidden layer followed by a scalar linear readout gives the cleanest structural case. Let
\[
h_t:\mathcal X\to\mathbb R^m,
\qquad
(T_t\beta)(x)=\beta^\top h_t(x),
\]
with representation span $\mathcal S_t=\Ran(T_t)$. Let $C_t=T_t^\ast T_t=\mathbb E[h_th_t^\top]$ and $K_t=T_tT_t^\ast$. Their positive spectra coincide.

For a neuron subset $I\subset\{1,\ldots,m\}$ let $T_{t,I}$ keep only those coordinates and let $P_t,P_{t,I}$ denote the corresponding prediction-space projectors.

\begin{proposition}[Exact current cost of neuron contraction]\label{prop:neuron-cost}
The increase in population risk of the optimal frozen linear readout caused by retaining only $I$ is
\[
\boxed{
D_t(I)=\norm{(P_t-P_{t,I})f^\star}^2\ge0.}
\]
If $\Ran(T_{t,I})=\Ran(T_t)$, then $D_t(I)=0$.
\end{proposition}

\begin{proof}
The optimal frozen-readout excess risk in a representation span $\mathcal S$ is
$\norm{(I-P_\mathcal S)f^\star}^2$. Since $\Ran(T_{t,I})\subseteq\Ran(T_t)$, the projectors are nested and Pythagoras yields the difference $\norm{(P_t-P_{t,I})f^\star}^2$.
\end{proof}

\begin{theorem}[A rank-$r$ representation contains an exact $r$-neuron basis]\label{thm:subset-basis}
If $\mathrm{rank}(T_t)=r< m$, there exists a subset $I$ of exactly $r$ neuron coordinates such that
\[
\Ran(T_{t,I})=\Ran(T_t).
\]
Therefore the current optimal frozen-readout risk is preserved exactly after deleting $m-r$ neurons.
\end{theorem}

\begin{proof}
The $m$ coordinate feature functions $h_{t,1},\ldots,h_{t,m}$ span the $r$-dimensional space $\Ran(T_t)$. Every finite spanning set contains a basis subset of size $r$.
\end{proof}

\begin{remark}[Why not immediately compress to one feature?]
For a single scalar target, if the population-optimal readout coefficient were known, one could collapse the \emph{current predictor} into the scalar feature $\beta_t^{\star\top}h_t$. That does not preserve future feature-learning capacity. The relevant object for training-time contraction is therefore not current predictive sufficiency alone but \emph{current sufficiency plus a bound on future discovery value}. This is the central difference between deployment compression and training-time contraction.
\end{remark}

\section{The onset of safe reduction in the exact quadratic phase model}
The endogenous spectral-geometry monograph provides a solvable nonlinear teacher--student model in which the pruning-onset question can be answered cleanly. The learned covariance $A_t$ has a transported random bulk with upper edge $b_+(t)$ satisfying
\[
b_+(t)\sim(4t)^{-1},
\]
while every fixed teacher eigenvalue $\theta_q$ has a unique dynamic BBP-type crossing time $\tau_{\theta_q,\gamma}$. Before that time the teacher mode has no separated outlier; afterward a unique outlier $z_q(t)>b_+(t)$ appears, converges to $\theta_q$, and its eigenvector overlap $\omega_q(t)$ with the teacher direction converges to one. The tangent prediction operator sends this to a spike $\kappa_q(t)=4z_q(t)$ above the prediction bulk edge $4b_+(t)$.

\begin{theorem}[Oracle phase-aware pruning onset in the quadratic model]\label{thm:bbp-onset}
Consider finitely many distinct positive teacher modes $\theta_1,\ldots,\theta_k$ satisfying the supercritical terminal conditions of the quadratic phase theorem. Define
\[
\tau_{\rm sep}=\max_{1\le q\le k}\tau_{\theta_q,\gamma}.
\]
Then:
\begin{enumerate}[label=(\roman*)]
\item No rule that identifies all teacher modes solely as detached covariance/prediction outliers can be correct for every $t<\tau_{\rm sep}$.
\item For every $\varepsilon>0$, there exists $T_\varepsilon>\tau_{\rm sep}$ such that for all $t\ge T_\varepsilon$,
\[
b_+(t)\le\varepsilon,
\qquad
1-\omega_q(t)\le\varepsilon\quad\forall q.
\]
\item Consequently, for sufficiently large $t$ there exists a deterministic threshold $c_t$ that separates all teacher-associated prediction spikes from the random bulk:
\[
4b_+(t)<c_t<\min_q4z_q(t).
\]
\end{enumerate}
Thus the solvable model has an endogenous discovery phase followed by an asymptotically identifiable compression phase.
\end{theorem}

\begin{proof}
Part (i) follows because at least one teacher mode is, by definition of $\tau_{\rm sep}$, still absorbed in the bulk before the last crossing. For (ii), the phase theorem gives $b_+(t)\to0$ and $\omega_q(t)\to1$ for each fixed $q$; finiteness of $k$ lets one choose a common $T_\varepsilon$. It also gives $z_q(t)\to\theta_q>0$. Hence for sufficiently large $t$, $4b_+(t)<2\min_q\theta_q$ while $4z_q(t)>2\min_q\theta_q$ after increasing $T_\varepsilon$ if needed, so a separating threshold exists, proving (iii).
\end{proof}

\subsection{Why premature rank-one contraction creates a discovery delay}
The preceding onset theorem concerns observability. In the rank-one teacher model one can go further and quantify the optimization price of contracting before a target direction has macroscopic overlap.

\begin{proposition}[Exact target-angle dynamics after rank-one contraction]\label{prop:rankone-angle}
Let the teacher be $A^\star=\theta uu^\top$ with $\theta>0$, and after a contraction let the student have one active column $w_s$, so $A_s=w_sw_s^\top$. Define
\[
\omega_s=\frac{(u^\top w_s)^2}{\norm{w_s}^2}\in[0,1].
\]
Under the exact population gradient flow
\[
\dot w_s=-2(w_sw_s^\top-A^\star)w_s,
\]
for every interval on which $0<\omega_s<1$,
\[
\boxed{
\frac{1-\omega_{t+h}}{\omega_{t+h}}
=
\frac{1-\omega_t}{\omega_t}e^{-4\theta h}.}
\]
Therefore the additional time required to reach squared teacher overlap at least $1-\varepsilon$ is
\[
\boxed{
T_\varepsilon(\omega_t)
=
\frac{1}{4\theta}
\left[
\log\frac{(1-\omega_t)(1-\varepsilon)}{\omega_t\varepsilon}
\right]_+.}
\]
\end{proposition}

\begin{proof}
Decompose $w=cu+z$ with $z\perp u$ and write $q=\norm w^2$. The gradient flow gives
\[
\dot c=-2(q-\theta)c,
\qquad
\dot z=-2qz.
\]
Hence
\[
\frac{d}{ds}\log\frac{\norm z^2}{c^2}=-4\theta.
\]
Since $(1-\omega)/\omega=\norm z^2/c^2$, integration yields the odds identity. Solving the inequality $\omega_{t+h}\ge1-\varepsilon$ gives the hitting time.
\end{proof}

\begin{corollary}[Discovery-delay separation before and after spectral resolution]\label{cor:discovery-delay}
At isotropic initialization, any fixed sample covariance eigenvector is Haar relative to an independent deterministic teacher direction. Its squared overlap satisfies $\omega_0=O_P(d^{-1})$. Consequently a blind rank-one contraction at initialization requires
\[
T_\varepsilon(\omega_0)
=
\frac{\log d}{4\theta}+O_P(1)
\]
additional population time to obtain constant teacher overlap. By contrast, if contraction is performed at a fixed supercritical checkpoint $t>\tau_{\theta,\gamma}$ and retains the teacher-associated outlier, Theorem~5.5 of the endogenous phase theory gives $\omega_t\to\omega(t)>0$, so $T_\varepsilon(\omega_t)=O(1)$. As $t\to\infty$, $\omega(t)\to1$ and the additional alignment delay vanishes.
\end{corollary}

\begin{proof}
For a Haar unit vector $v$ independent of $u$, $d|u^\top v|^2$ is tight, so $\log(1/\omega_0)=\log d+O_P(1)$. Apply Proposition~\ref{prop:rankone-angle}. In the supercritical regime the outlier-overlap theorem supplies a nonzero deterministic limit, and its long-time limit is one.
\end{proof}

This result sharpens the meaning of ``do not prune too early.'' Premature contraction need not destroy the teacher forever: infinitesimal overlap can eventually be amplified. What it does is convert an order-one resolved feature into a dimension-dependent discovery delay. The BBP crossing is therefore a natural observability boundary for avoiding that delay, not a mystical universal threshold.

The theorem does \emph{not} claim that a BBP threshold is a universal pruning rule for arbitrary networks. It establishes the exact mechanism in one nonlinear model and provides the theory-derived reason for a warm-up/separation phase.
